\IfFileExists{stacks-project.cls}{%
\documentclass{stacks-project}
}{%
\documentclass{amsart}
}

% For dealing with references we use the comment environment
\usepackage{verbatim}
\newenvironment{reference}{\comment}{\endcomment}
%\newenvironment{reference}{}{}
\newenvironment{slogan}{\comment}{\endcomment}
\newenvironment{history}{\comment}{\endcomment}

% For commutative diagrams we use Xy-pic
\usepackage[all]{xy}

% We use 2cell for 2-commutative diagrams.
\xyoption{2cell}
\UseAllTwocells

% We use multicol for the list of chapters between chapters
\usepackage{multicol}

% This is generally recommended for better output
\usepackage{lmodern}
\usepackage[T1]{fontenc}

% For cross-file-references
\usepackage{xr-hyper}

% Package for hypertext links:
\usepackage{hyperref}

% For any local file, say "hello.tex" you want to link to please
% use \externaldocument[hello-]{hello}
\externaldocument[introduction-]{introduction}
\externaldocument[conventions-]{conventions}
\externaldocument[sets-]{sets}
\externaldocument[categories-]{categories}
\externaldocument[topology-]{topology}
\externaldocument[sheaves-]{sheaves}
\externaldocument[sites-]{sites}
\externaldocument[stacks-]{stacks}
\externaldocument[fields-]{fields}
\externaldocument[algebra-]{algebra}
\externaldocument[brauer-]{brauer}
\externaldocument[homology-]{homology}
\externaldocument[derived-]{derived}
\externaldocument[simplicial-]{simplicial}
\externaldocument[more-algebra-]{more-algebra}
\externaldocument[smoothing-]{smoothing}
\externaldocument[modules-]{modules}
\externaldocument[sites-modules-]{sites-modules}
\externaldocument[injectives-]{injectives}
\externaldocument[cohomology-]{cohomology}
\externaldocument[sites-cohomology-]{sites-cohomology}
\externaldocument[dga-]{dga}
\externaldocument[dpa-]{dpa}
\externaldocument[sdga-]{sdga}
\externaldocument[hypercovering-]{hypercovering}
\externaldocument[schemes-]{schemes}
\externaldocument[constructions-]{constructions}
\externaldocument[properties-]{properties}
\externaldocument[morphisms-]{morphisms}
\externaldocument[coherent-]{coherent}
\externaldocument[divisors-]{divisors}
\externaldocument[limits-]{limits}
\externaldocument[varieties-]{varieties}
\externaldocument[topologies-]{topologies}
\externaldocument[descent-]{descent}
\externaldocument[perfect-]{perfect}
\externaldocument[more-morphisms-]{more-morphisms}
\externaldocument[flat-]{flat}
\externaldocument[groupoids-]{groupoids}
\externaldocument[more-groupoids-]{more-groupoids}
\externaldocument[etale-]{etale}
\externaldocument[chow-]{chow}
\externaldocument[intersection-]{intersection}
\externaldocument[pic-]{pic}
\externaldocument[weil-]{weil}
\externaldocument[adequate-]{adequate}
\externaldocument[dualizing-]{dualizing}
\externaldocument[duality-]{duality}
\externaldocument[discriminant-]{discriminant}
\externaldocument[derham-]{derham}
\externaldocument[local-cohomology-]{local-cohomology}
\externaldocument[algebraization-]{algebraization}
\externaldocument[curves-]{curves}
\externaldocument[resolve-]{resolve}
\externaldocument[models-]{models}
\externaldocument[functors-]{functors}
\externaldocument[equiv-]{equiv}
\externaldocument[pione-]{pione}
\externaldocument[etale-cohomology-]{etale-cohomology}
\externaldocument[proetale-]{proetale}
\externaldocument[relative-cycles-]{relative-cycles}
\externaldocument[more-etale-]{more-etale}
\externaldocument[trace-]{trace}
\externaldocument[crystalline-]{crystalline}
\externaldocument[spaces-]{spaces}
\externaldocument[spaces-properties-]{spaces-properties}
\externaldocument[spaces-morphisms-]{spaces-morphisms}
\externaldocument[decent-spaces-]{decent-spaces}
\externaldocument[spaces-cohomology-]{spaces-cohomology}
\externaldocument[spaces-limits-]{spaces-limits}
\externaldocument[spaces-divisors-]{spaces-divisors}
\externaldocument[spaces-over-fields-]{spaces-over-fields}
\externaldocument[spaces-topologies-]{spaces-topologies}
\externaldocument[spaces-descent-]{spaces-descent}
\externaldocument[spaces-perfect-]{spaces-perfect}
\externaldocument[spaces-more-morphisms-]{spaces-more-morphisms}
\externaldocument[spaces-flat-]{spaces-flat}
\externaldocument[spaces-groupoids-]{spaces-groupoids}
\externaldocument[spaces-more-groupoids-]{spaces-more-groupoids}
\externaldocument[bootstrap-]{bootstrap}
\externaldocument[spaces-pushouts-]{spaces-pushouts}
\externaldocument[spaces-chow-]{spaces-chow}
\externaldocument[groupoids-quotients-]{groupoids-quotients}
\externaldocument[spaces-more-cohomology-]{spaces-more-cohomology}
\externaldocument[spaces-simplicial-]{spaces-simplicial}
\externaldocument[spaces-duality-]{spaces-duality}
\externaldocument[formal-spaces-]{formal-spaces}
\externaldocument[restricted-]{restricted}
\externaldocument[spaces-resolve-]{spaces-resolve}
\externaldocument[formal-defos-]{formal-defos}
\externaldocument[defos-]{defos}
\externaldocument[cotangent-]{cotangent}
\externaldocument[examples-defos-]{examples-defos}
\externaldocument[algebraic-]{algebraic}
\externaldocument[examples-stacks-]{examples-stacks}
\externaldocument[stacks-sheaves-]{stacks-sheaves}
\externaldocument[criteria-]{criteria}
\externaldocument[artin-]{artin}
\externaldocument[quot-]{quot}
\externaldocument[stacks-properties-]{stacks-properties}
\externaldocument[stacks-morphisms-]{stacks-morphisms}
\externaldocument[stacks-limits-]{stacks-limits}
\externaldocument[stacks-cohomology-]{stacks-cohomology}
\externaldocument[stacks-perfect-]{stacks-perfect}
\externaldocument[stacks-introduction-]{stacks-introduction}
\externaldocument[stacks-more-morphisms-]{stacks-more-morphisms}
\externaldocument[stacks-geometry-]{stacks-geometry}
\externaldocument[moduli-]{moduli}
\externaldocument[moduli-curves-]{moduli-curves}
\externaldocument[examples-]{examples}
\externaldocument[exercises-]{exercises}
\externaldocument[guide-]{guide}
\externaldocument[desirables-]{desirables}
\externaldocument[coding-]{coding}
\externaldocument[obsolete-]{obsolete}
\externaldocument[fdl-]{fdl}
\externaldocument[index-]{index}

% Theorem environments.
%
\theoremstyle{plain}
\newtheorem{theorem}[subsection]{Theorem}
\newtheorem{proposition}[subsection]{Proposition}
\newtheorem{lemma}[subsection]{Lemma}

\theoremstyle{definition}
\newtheorem{definition}[subsection]{Definition}
\newtheorem{example}[subsection]{Example}
\newtheorem{exercise}[subsection]{Exercise}
\newtheorem{situation}[subsection]{Situation}

\theoremstyle{remark}
\newtheorem{remark}[subsection]{Remark}
\newtheorem{remarks}[subsection]{Remarks}

\numberwithin{equation}{subsection}

% Macros
%
\def\lim{\mathop{\mathrm{lim}}\nolimits}
\def\colim{\mathop{\mathrm{colim}}\nolimits}
\def\Spec{\mathop{\mathrm{Spec}}}
\def\Hom{\mathop{\mathrm{Hom}}\nolimits}
\def\Ext{\mathop{\mathrm{Ext}}\nolimits}
\def\SheafHom{\mathop{\mathcal{H}\!\mathit{om}}\nolimits}
\def\SheafExt{\mathop{\mathcal{E}\!\mathit{xt}}\nolimits}
\def\Sch{\mathit{Sch}}
\def\Mor{\mathop{\mathrm{Mor}}\nolimits}
\def\Ob{\mathop{\mathrm{Ob}}\nolimits}
\def\Sh{\mathop{\mathit{Sh}}\nolimits}
\def\NL{\mathop{N\!L}\nolimits}
\def\CH{\mathop{\mathrm{CH}}\nolimits}
\def\proetale{{pro\text{-}\acute{e}tale}}
\def\etale{{\acute{e}tale}}
\def\QCoh{\mathit{QCoh}}
\def\Ker{\mathop{\mathrm{Ker}}}
\def\Im{\mathop{\mathrm{Im}}}
\def\Coker{\mathop{\mathrm{Coker}}}
\def\Coim{\mathop{\mathrm{Coim}}}

% Boxtimes
%
\DeclareMathSymbol{\boxtimes}{\mathbin}{AMSa}{"02}

%
% Macros for moduli stacks/spaces
%
\def\QCohstack{\mathcal{QC}\!\mathit{oh}}
\def\Cohstack{\mathcal{C}\!\mathit{oh}}
\def\Spacesstack{\mathcal{S}\!\mathit{paces}}
\def\Quotfunctor{\mathrm{Quot}}
\def\Hilbfunctor{\mathrm{Hilb}}
\def\Curvesstack{\mathcal{C}\!\mathit{urves}}
\def\Polarizedstack{\mathcal{P}\!\mathit{olarized}}
\def\Complexesstack{\mathcal{C}\!\mathit{omplexes}}
% \Pic is the operator that assigns to X its picard group, usage \Pic(X)
% \Picardstack_{X/B} denotes the Picard stack of X over B
% \Picardfunctor_{X/B} denotes the Picard functor of X over B
\def\Pic{\mathop{\mathrm{Pic}}\nolimits}
\def\Picardstack{\mathcal{P}\!\mathit{ic}}
\def\Picardfunctor{\mathrm{Pic}}
\def\Deformationcategory{\mathcal{D}\!\mathit{ef}}
\usepackage{cleveref}
\crefname{originallemma}{original lemma}{original lemmas}
\Crefname{originallemma}{Original lemma}{Original lemmas}
\makeatletter
\newlabel{derived-lemma-two-ss-complex-functor}{{015J}{source}{}{sourcepointer-1}{}}
\newlabel{derived-lemma-two-ss-complex-functor@cref}{{[originallemma][0][]015J}{[1][0][]source}{}{}{}}
\makeatother
\crefname{originallemma}{original lemma}{original lemmas}
\Crefname{originallemma}{Original lemma}{Original lemmas}
\makeatletter
\newlabel{derived-lemma-unbounded-right-derived}{{07K7}{source}{}{sourcepointer-2}{}}
\newlabel{derived-lemma-unbounded-right-derived@cref}{{[originallemma][0][]07K7}{[1][0][]source}{}{}{}}
\makeatother
\crefname{originallemma}{original lemma}{original lemmas}
\Crefname{originallemma}{Original lemma}{Original lemmas}
\makeatletter
\newlabel{etale-cohomology-lemma-cohomological-dimension-proper}{{095U}{source}{}{sourcepointer-3}{}}
\newlabel{etale-cohomology-lemma-cohomological-dimension-proper@cref}{{[originallemma][0][]095U}{[1][0][]source}{}{}{}}
\makeatother
\crefname{originallemma}{original lemma}{original lemmas}
\Crefname{originallemma}{Original lemma}{Original lemmas}
\makeatletter
\newlabel{etale-cohomology-lemma-constructible-abelian}{{03RZ}{source}{}{sourcepointer-4}{}}
\newlabel{etale-cohomology-lemma-constructible-abelian@cref}{{[originallemma][0][]03RZ}{[1][0][]source}{}{}{}}
\makeatother
\crefname{originallemma}{original lemma}{original lemmas}
\Crefname{originallemma}{Original lemma}{Original lemmas}
\makeatletter
\newlabel{etale-cohomology-lemma-projection-formula-proper-mod-n}{{0F0G}{source}{}{sourcepointer-5}{}}
\newlabel{etale-cohomology-lemma-projection-formula-proper-mod-n@cref}{{[originallemma][0][]0F0G}{[1][0][]source}{}{}{}}
\makeatother
\crefname{originallemma}{original lemma}{original lemmas}
\Crefname{originallemma}{Original lemma}{Original lemmas}
\makeatletter
\newlabel{etale-cohomology-lemma-tensor-constructible}{{0GKB}{source}{}{sourcepointer-6}{}}
\newlabel{etale-cohomology-lemma-tensor-constructible@cref}{{[originallemma][0][]0GKB}{[1][0][]source}{}{}{}}
\makeatother
\crefname{originallemma}{original lemma}{original lemmas}
\Crefname{originallemma}{Original lemma}{Original lemmas}
\makeatletter
\newlabel{lemma-proper-smooth-family-curves}{{0GKV}{source}{}{sourcepointer-7}{}}
\newlabel{lemma-proper-smooth-family-curves@cref}{{[originallemma][0][]0GKV}{[1][0][]source}{}{}{}}
\makeatother
\crefname{originallemma}{original lemma}{original lemmas}
\Crefname{originallemma}{Original lemma}{Original lemmas}
\makeatletter
\newlabel{sites-cohomology-lemma-modules-abelian-unbounded}{{0D6J}{source}{}{sourcepointer-8}{}}
\newlabel{sites-cohomology-lemma-modules-abelian-unbounded@cref}{{[originallemma][0][]0D6J}{[1][0][]source}{}{}{}}
\makeatother
\title{Constructibility of R\textasciicircum{}qf\_* of a torsion finite module along a proper smooth family of curves with geometrically connected fibres}
\author{The Stacks Project Authors}
\date{stacks-project:more-etale.tex:local-snapshot-20261012; GNU FDL 1.2 or later}
\begin{document}
\maketitle
\setcounter{subsection}{1}
\noindent\textit{Editorial note.} All delivered mathematical text is the author's own, copied from the pinned Stacks Project chapter file more-etale.tex; no statement or proof step was rewritten, reordered or removed. The counted span is the author environment and its own proof environment. References inside the span are handed to the assembler as navigation stubs that carry the canonical tag of the cited tag only; whether each of them is load-bearing is recorded on the stub itself.\par\medskip
\noindent\textit{Allowed standard prerequisites (editorial).}\begin{itemize}
\item Constructible abelian sheaves closed under extensions (Etale Cohomology, Lemma etale-cohomology-lemma-constructible-abelian, canonical tag 03RZ). cited at proof lines 4604-4606. Runs the induction on the torsion degree: the long exact sequence sandwiches R\textasciicircum{}qf\_*M between constructible terms once the prime case is known.
\item Projection formula for proper pushforward of mod-n sheaves (Etale Cohomology, Lemma etale-cohomology-lemma-projection-formula-proper-mod-n, canonical tag 0F0G). cited at proof lines 4649-4656. Trades Rf\_*Lambda for the derived tensor of Rf\_*Z/lZ with Lambda, connecting the free-module case to the constant sheaf case.
\item Constant sheaf case (Lemma lemma-proper-smooth-family-curves, canonical tag 0GKV). cited at proof line 4667. The terminal base case: constructibility of R\textasciicircum{}qf\_*Z/lZ for the proper smooth family of curves.
\end{itemize}\medskip
\ifdefined\hypertarget\hypertarget{sourcepointer-1}{}\fi
\noindent\textit{Source reference.} Original lemma 015J: sources/stacks-chow/tags.txt, source lines 1510--1510. Cited at proof lines 4642-4644: the first spectral sequence reduces constructibility of the cohomology sheaves of the truncated free complex to the free terms.\par\medskip
\ifdefined\hypertarget\hypertarget{sourcepointer-2}{}\fi
\noindent\textit{Source reference.} Original lemma 07K7: sources/stacks-chow/tags.txt, source lines 9647--9647. Cited at proof lines 4629-4634: with the finite cohomological dimension lemma it truncates the free resolution to a bounded complex.\par\medskip
\ifdefined\hypertarget\hypertarget{sourcepointer-3}{}\fi
\noindent\textit{Source reference.} Original lemma 095U: sources/stacks-chow/tags.txt, source lines 11597--11597. Cited at proof lines 4629-4632: f\_* has finite cohomological dimension on sheaves of Lambda-modules, enabling the truncation.\par\medskip
\ifdefined\hypertarget\hypertarget{sourcepointer-4}{}\fi
\noindent\textit{Source reference.} Original lemma 03RZ: sources/stacks-chow/tags.txt, source lines 4948--4948. Cited at proof lines 4604-4606: the outer terms of the long exact sequence are constructible, winning the induction on n.\par\medskip
\ifdefined\hypertarget\hypertarget{sourcepointer-5}{}\fi
\noindent\textit{Source reference.} Original lemma 0F0G: sources/stacks-chow/tags.txt, source lines 18761--18761. Cited at proof lines 4649-4656: (Rf\_*Z/lZ) tensor\textasciicircum{}L Lambda = Rf\_*Lambda, the key trade to the constant sheaf case.\par\medskip
\ifdefined\hypertarget\hypertarget{sourcepointer-6}{}\fi
\noindent\textit{Source reference.} Original lemma 0GKB: sources/stacks-chow/tags.txt, source lines 20681--20681. Cited at proof lines 4664-4666: tensoring the constructible R\textasciicircum{}qf\_*Z/lZ with Lambda stays constructible.\par\medskip
\ifdefined\hypertarget\hypertarget{sourcepointer-7}{}\fi
\noindent\textit{Source reference.} Original lemma 0GKV: sources/stacks-chow/tags.txt, source lines 20700--20700. Cited at proof line 4667: the terminal base case, constructibility of R\textasciicircum{}qf\_*Z/lZ.\par\medskip
\ifdefined\hypertarget\hypertarget{sourcepointer-8}{}\fi
\noindent\textit{Source reference.} Original lemma 0D6J: sources/stacks-chow/tags.txt, source lines 16524--16524. Cited at proof lines 4611-4616: R\textasciicircum{}qf\_*M does not change when M is viewed as a sheaf of Lambda/lLambda-modules.\par\medskip
\begin{lemma}
\label{lemma-proper-smooth-family-curves-modules}
Let $f : X \to S$ be a proper smooth morphism of schemes with
geometrically connected fibres of dimension $1$. Let $\Lambda$
be a Noetherian ring. Let $M$ be a finite $\Lambda$-module
annihilated by an integer $n > 0$. Then $R^qf_*\underline{M}$
is a constructible sheaf of $\Lambda$-modules on $S$.
\end{lemma}

\begin{proof}
If $n = \ell n'$ for some prime number $\ell$, then
we get a short exact sequence $0 \to M[\ell] \to M \to M' \to 0$
of finite $\Lambda$-modules and $M'$ is annihilated by $n'$.
This produces a corresponding short exact sequence of constant
sheaves, which in turn gives rise to an exact sequence
$$
R^{q - 1}f_*\underline{M'} \to
R^qf_*\underline{M[n]} \to
R^qf_*\underline{M} \to
R^qf_*\underline{M'} \to
R^{q + 1}f_*\underline{M[n]}
$$
Thus, if we can show the result in case $M$ is annihilated by
a prime number, then by induction on $n$ we win by
\'Etale Cohomology, Lemma \ref{etale-cohomology-lemma-constructible-abelian}.

\medskip\noindent
Let $\ell$ be a prime number such that $\ell$ annihilates $M$.
Then we can replace $\Lambda$ by the $\mathbf{F}_\ell$-algebra
$\Lambda/\ell \Lambda$. Namely, the sheaf $R^qf_*\underline{M}$
where $\underline{M}$ is viewed as a sheaf of $\Lambda$-modules
is the same as the sheaf $R^qf_*\underline{M}$ computed by
viewing $\underline{M}$ as a sheaf of $\Lambda/\ell \Lambda$-modules, see
Cohomology on Sites, Lemma
\ref{sites-cohomology-lemma-modules-abelian-unbounded}.

\medskip\noindent
Assume $\ell$ be a prime number such that $\ell$ annihilates $M$ and $\Lambda$.
Let us reduce to the case where $M$ is a finite free $\Lambda$-module.
Namely, choose a resolution
$$
\ldots \to
\Lambda^{\oplus m_2} \to
\Lambda^{\oplus m_1} \to
\Lambda^{\oplus m_0} \to
M \to 0
$$
Recall that $f_*$ has finite cohomological dimension on sheaves of
$\Lambda$-modules, see
\'Etale Cohomology, Lemma
\ref{etale-cohomology-lemma-cohomological-dimension-proper}
and
Derived Categories, Lemma \ref{derived-lemma-unbounded-right-derived}.
Thus we see that $R^qf_*\underline{M}$ is the $q$th cohomology
sheaf of the object
$$
Rf_*(\underline{\Lambda^{\oplus m_a}} \to \ldots \to
\underline{\Lambda^{\oplus m_0}})
$$
in $D(S_\etale, \Lambda)$ for some integer $a$ large enough.
Using the first spectral sequence of
Derived Categories, Lemma \ref{derived-lemma-two-ss-complex-functor}
(or alternatively using an argument with truncations)
we conclude that it suffices to prove that $R^qf_*\underline{\Lambda})$
is constructible.

\medskip\noindent
At this point we can finally use that
$$
(Rf_*\underline{\mathbf{Z}/\ell\mathbf{Z}})
\otimes_{\mathbf{Z}/\ell\mathbf{Z}}^\mathbf{L} \underline{\Lambda} =
Rf_*\underline{\Lambda}
$$
by \'Etale Cohomology, Lemma
\ref{etale-cohomology-lemma-projection-formula-proper-mod-n}.
Since any module over the field $\mathbf{Z}/\ell\mathbf{Z}$
is flat we obtain
$$
(R^qf_*\underline{\mathbf{Z}/\ell\mathbf{Z}})
\otimes_{\mathbf{Z}/\ell\mathbf{Z}} \underline{\Lambda} =
R^qf_*\underline{\Lambda}
$$
Hence it suffices to prove the result for
$R^qf_*\underline{\mathbf{Z}/\ell\mathbf{Z}}$
by \'Etale Cohomology, Lemma \ref{etale-cohomology-lemma-tensor-constructible}.
This case is Lemma \ref{lemma-proper-smooth-family-curves}.
\end{proof}

\end{document}
